\documentclass{article}
\usepackage[utf8]{inputenc}
\usepackage[french]{babel}
\usepackage{amsmath}
\usepackage{siunitx}
\usepackage[version=4]{mhchem}
\usepackage{tkz-euclide}
\usepackage{pgfplots}
%\usepackage{graphicx}

\usetikzlibrary{intersections,angles,quotes}
\usepgfplotslibrary{fillbetween}

\begin{document}

\section{Le cercle trigonométrique}

Le cercle trigonométrique est un outil fondamental pour comprendre les fonctions trigonométriques.
Il s'agit d'un cercle de rayon 1 centré à l'origine d'un repère orthonormé.

\subsection{Définition}

Un cercle trigonométrique est un cercle~:
\begin{itemize}
    \item de centre $O$ (l'origine du repère)~;
    \item de rayon $1$~;
    \item orienté dans le sens inverse des aiguilles d'une montre (sens trigonométrique).
\end{itemize}

\subsection{Repérage d'un point sur le cercle}

Tout point $M$ du cercle trigonométrique peut être repéré par un angle $\alpha$ (en radians ou en degrés) formé entre l'axe des abscisses et la demi-droite  $[O, M)$.

\begin{figure}[h]
    \centering
\begin{tikzpicture}[scale=5]
    \clip (-0.25,-0.25) rectangle (1.5,1.25);
    \coordinate (O) at (0,0);
    \node[below left]  at (O) {$O$};
    %grid lines
    \draw[
            step=.5cm,
            gray,
            very thin
            ]
            (-1.2,-1.2) grid (1.4,1.2);
    %\filldraw[
    %            fill=blue!20,
    %            draw=red!50
    %            ]
    %            (0,0) -- (3mm,0mm)
    %arc [
    %        start angle=0,
    %        end angle=30,
    %        radius=3mm]
    %        -- cycle;
    %axes
    \draw[->] (-1.2,0) -- (1.5,0) coordinate (x axis);
    \draw[->] (0,-1.2) -- (0,1.2) coordinate (y axis);

    %axes label
    %\node [right]at (1.5,0)(x){$x$};
    %\node [above]at (0,1.5){$y$};
    %circle
    \draw (0,0) circle [
                        radius=1cm
                        ];
    % Point M on the circle at angle alpha (30 degrees)
    \coordinate (M) at (30:1cm);
    \filldraw (M) circle (0.4pt) node[above left, inner sep=1pt] {$M$};
    %triangle height
    \draw[
        very thick,
        orange
        ]
        (30:1cm) -- node[
                        left=1pt,
                        fill=white
                        ]
                        {$\sin \alpha$} (30:1cm |- x axis);
    %triangle base
    \draw[
            very thick,
            blue
            ]
            (30:1cm |- x axis) -- node[
                                        above=1pt,
                                        xshift=5pt,
                                        fill=white
                                        ]
                                        {$\cos \alpha$} (0,0);
    \coordinate (x) at (30:1cm |- x axis);
    \node[below]  at (x) {$x$};
    \coordinate (y) at (0,.5);
    \node[left]  at (y) {$y$};
   %intersection
    \path [name path=upward line] (1,0) -- (1,1);
    \path [name path=sloped line] (0,0) -- (30:1.5cm);
    \draw [name intersections={of=upward line and sloped line, by=t}]
            [very thick,red]
            (1,0) -- node [right=1pt,fill=white]
            {$\displaystyle \tan \alpha $} (t);
    \draw (0,0) -- (t);

    %x-ticks
    \foreach \x/\xtext in {-1,
                            -0.5/-\frac{1}{2},
                                            1}
    \draw (\x cm,1pt) -- (\x cm,-1pt) node[
                                            anchor=north,
                                            fill=white
                                            ]
                                            {$\xtext$};
    %y-ticks
    \foreach \y/\ytext in {-1, -0.5/-\frac{1}{2}, 0.5/, 1}
    \draw (1pt,\y cm) -- (-1pt,\y cm)
        \ifx\ytext\empty\else
            node[anchor=east, fill=white] {$\ytext$}
        \fi;

    %arc angle
    \draw (x) coordinate (A)--
        (0,0) coordinate (B)--
        (t) coordinate (C)
    pic [
            draw,
            red,
            "$\alpha$",
            angle radius=9mm
        ] {angle};
    
        \draw[thick] (O) -- (x) node[midway,below] {$a$};
        \draw[thick] (x) -- (M) node[midway,right] {$b$};
        \draw[thick] (O) -- (M) node[midway,above right, yshift=2pt] {$c$};
\end{tikzpicture}
    \label{fig:trigo}
    \caption{Cercle trigonométrique}
\end{figure}

\subsection{Coordonnées d'un point}
Pour un point $M$ sur le cercle trigonométrique, repéré par un angle $\alpha$, ses coordonnées $(x, y)$ sont données par~:
\[
x = \cos(\alpha) \quad \text{et} \quad y = \sin(\alpha)
\]

\section{Rappel du théorème de Pythagore}

Le théorème de Pythagore s'applique au triangle rectangle formé par les côtés \(a\), \(b\) et \(c\) (où \(c\) est l'hypoténuse) de la figure~\ref{fig:trigo}.
Il s'énonce ainsi~:
\[
c^2 = a^2 + b^2
\]

\subsection{Application au cercle trigonométrique}
Dans le cercle trigonométrique, pour un point $M (x,y)$ sur le cercle, ses coordonnées vérifient~:
\[
x^2 + y^2 = 1
\]
Cela découle directement du théorème de Pythagore appliqué au triangle rectangle de la figure~\ref{fig:trigo}.

\section{Fonctions trigonométriques de base}

Les fonctions trigonométriques permettent de relier les angles d'un triangle rectangle à ses côtés.

\subsection{Fonctions sinus, cosinus et tangente}

\begin{itemize}
    \item \(\sin(\alpha) = \frac{a}{c}\)
    \item \(\cos(\alpha) = \frac{b}{c}\)
    \item \(\tan(\alpha) = \frac{a}{b} = \frac{\sin(\alpha)}{\cos(\alpha)}\)
\end{itemize}

\subsection{Fonctions réciproques}
Les fonctions réciproques permettent de retrouver l'angle $\alpha$ à partir de la valeur de la fonction trigonométrique~:
\begin{itemize}
    \item $\arcsin(x)$~: angle dont le sinus est  x
    \item $\arccos(x)$~: angle dont le cosinus est  x
    \item $\arctan(x)$~: angle dont la tangente est  x
\end{itemize}

\smallskip
\textbf{Exemple}~:
Si $\sin(\alpha) = 0{,}5$, alors $\alpha = \arcsin(0{,}5) = \frac{\pi}{6}$(ou \SI{30}{\degree}).

\section{Identités trigonométriques remarquables}

\subsection{Identité fondamentale}  

Pour tout angle  $\alpha$~:  
\[
\sin^2(\alpha) + \cos^2(\alpha) = 1  
\]

(découle directement du théorème de Pythagore.)

\subsection{Formules d'addition}  
Pour tous angles  $a$  et  $b$~:  
\begin{align*}  
    \cos(a + b) &= \cos(a)\cos(b) - \sin(a)\sin(b) \\
    \cos(a - b) &= \cos(a)\cos(b) + \sin(a)\sin(b) \\
    \sin(a + b) &= \sin(a)\cos(b) + \cos(a)\sin(b) \\
    \sin(a - b) &= \sin(a)\cos(b) - \cos(a)\sin(b)
\end{align*}

(démontrées par projection des points sur le cercle trigonométrique.)

\subsection{Formules de duplication}  
Pour tout angle $\alpha$~:  
\begin{align*}  
    \sin(2\alpha) &= 2\sin(\alpha)\cos(\alpha) \\
    \cos(2\alpha) &= \cos^2(\alpha) - \sin^2(\alpha) \\
                  &= 2\cos^2(\alpha) - 1 \\
                  & = 1 - 2\sin^2(\alpha)
\end{align*}

(découlent des formules d'addition avec $a = b = \alpha$.)

\section{Quelques valeurs remarquables}

\begin{table}[h]
\centering
\renewcommand{\arraystretch}{1.3} % Aère les lignes verticalement
\begin{tabular}{c|cccc}
\toprule
$\alpha$ & $\sin\alpha$ & $\cos\alpha$ & $\tan\alpha$ & $\cot\alpha$ \\
\midrule
$0$ & $0$ & $1$ & $0$ & $\infty$ \\[6pt] % Espace vertical supplémentaire
$\frac{\pi}{6}$ ($30^\circ$) & $\frac{1}{2}$ & $\frac{\sqrt{3}}{2}$ & $\frac{\sqrt{3}}{3}$ & $\sqrt{3}$ \\[6pt]
$\frac{\pi}{4}$ ($45^\circ$) & $\frac{\sqrt{2}}{2}$ & $\frac{\sqrt{2}}{2}$ & $1$ & $1$ \\[6pt]
$\frac{\pi}{3}$ ($60^\circ$) & $\frac{\sqrt{3}}{2}$ & $\frac{1}{2}$ & $\sqrt{3}$ & $\frac{\sqrt{3}}{3}$ \\[6pt]
$\frac{\pi}{2}$ ($90^\circ$) & $1$ & $0$ & $\infty$ & $0$ \\
\bottomrule
\end{tabular}
\caption{Valeurs des fonctions trigonométriques pour des angles courants.}
\end{table}

(il est possible de retrouver certaines de ces valeurs juste en regardant le cercle trigonométrique.)

\end{document}
